\label{sec:conclusion}

Section~203 of the Voting Rights Act is the most overlooked provision
in redistricting analysis. Most redistricting commentary focuses on
Section~2 (majority-minority districts) and, historically, Section~5
(preclearance, now effectively dormant). But Section~203 creates
independent legal obligations that affect every redistricting cycle in
the 331 jurisdictions covered by the 2023 Census Bureau determination:
those jurisdictions must provide election materials in the applicable
language minority language regardless of how district lines are drawn.

Redistricting affects Section~203 not by changing which jurisdictions are
covered---coverage is determined by Census ACS data and is independent of
district boundaries---but by affecting the administrative ease of compliance
and the community integrity of language minority groups within the
redistricting unit. When a Section~203-covered county is split between
two congressional districts, the county election office faces increased
administrative burden. When the language minority community itself is split
across district lines, the political representation and advocacy capacity
of that community is diluted even as its Section~203 rights remain intact.

The \textsc{bisect} algorithmic redistricting pipeline addresses both
concerns. The county-weight mode (\texttt{--weights-override county})
minimizes Section~203 county splits by penalizing cross-county boundary
cuts in the METIS graph partitioning. For Texas, this produces a 35\%
reduction in Section~203 covered county splits relative to the enacted
2022 congressional plan; for California, a 9\% reduction. For the
Vietnamese community in Santa Clara County, \textsc{bisect} preserves
3 of 4 Vietnamese-majority census tracts within a single district,
compared with 1 of 4 under the enacted plan.

For jurisdictions with both Section~203 coverage and Section~2
majority-minority district requirements, \textsc{bisect} applies both
in sequence: the VRA mode (\texttt{--weights-override vra-aligned})
satisfies Section~2's majority-minority district requirement, and within
that constrained solution space, the county-weight optimization minimizes
Section~203 administrative burden. The 23 counties nationally with both
requirements present the most demanding redistricting context; the
\textsc{bisect} dual-optimization approach provides a principled,
verifiable resolution.

Expert witnesses and redistricting practitioners in states with significant
Section~203 coverage should incorporate both metrics in their analysis:
county split rates (for administrative burden) and community integrity
scores (for political representation). Both are straightforward to compute
from \textsc{bisect} output data, and both are relevant to the comprehensive
VRA compliance analysis that courts and redistricting commissions require.

\bigskip
\noindent\textit{The author thanks the redistricting research team for
data access and algorithmic development. All errors are the author's own.}
